% !TeX root = ../main.tex

\chapter{非加性最短路径问题及其算法}

本章主要围绕非加性最短路径问题的基本概念、典型特征及其求解方法展开论述.首先介绍图的基本定义及常见存储方式，为后续算法描述统一符号体系和数据结构基础；其次回顾经典最短路径问题的定义及其基本求解算法，说明传统最短路径模型成立所依赖的加性假设；然后在此基础上给出非加性最短路径问题的形式化定义、主要类别以及路径依赖型问题的经典求解思路，并分析传统算法在该类问题中失效的根本原因；接着介绍线图理论的基本概念及其在路径依赖建模中的作用；最后对 CUDA 并行框架进行简要说明，为下一章的算法设计与并行实现奠定基础.

\section{预备知识}

\subsection{图的定义}

图（Graph）是用于描述对象之间连接关系的一类基本数学结构.本文将图记作 $G=(V,E)$，其中 $V$ 表示顶点集合，$E$ 表示边集合；若图中每条边均带有相应代价，则可进一步记为 $G=(V,E,w)$，其中 $w:E\rightarrow\mathbb{R}_{\ge 0}$ 表示边权函数.通常将图的顶点数记为 $|V|=n$，边数记为 $|E|=m$.

为了便于说明后续的相关符号，图~\ref{fig:graph_example} 给出了一个简单的有向带权网络示例.在图论中，根据边是否具有方向，图可分为有向图与无向图.以图~\ref{fig:graph_example} 为例，边 $(A,B)$ 表示允许从顶点 $A$ 到达顶点 $B$，但并不意味着一定允许从 $B$ 返回 $A$.对于有向图，通常分别定义顶点的出度与入度，其中顶点 $v$ 的出度记为 $\deg^+(v)$，入度记为 $\deg^-(v)$；而在无向图中，由于边不区分方向，通常仅讨论顶点的度数 $\deg(v)$.

\begin{figure}[H]
  \centering
  \begin{tikzpicture}[
    >=Latex,
    node distance=22mm,
    every node/.style={circle,draw,minimum size=7mm,inner sep=0pt},
    every edge/.style={draw,->,line width=0.4pt}
  ]
    \node (A) {A};
    \node (B) [right=of A] {B};
    \node (C) [below=of A] {C};
    \node (D) [right=of C] {D};

    \path (A) edge node[above] {$3$} (B)
          (A) edge node[left]  {$5$} (C)
          (B) edge node[right] {$2$} (D)
          (C) edge node[below] {$4$} (D)
          (C) edge[bend left=15] node[above] {$1$} (B);
  \end{tikzpicture}
  \caption{网络图示例}
  \label{fig:graph_example}
\end{figure}


若边上附带距离、时间、能耗、风险等代价信息，则称该图为带权图；反之，若边仅表示连接关系而不包含额外代价，则称为无权图.由于本文讨论的最短路径问题本质上是以路径代价最小化为目标，因此后文默认研究对象均为带权图.

\subsection{图的存储格式}

图的存储方式会直接影响后续算法的访问效率与空间开销.若采用邻接矩阵进行存储，对于含有 $n$ 个顶点的图，需要 $O(n^2)$ 的存储空间.当图较为稀疏，即满足 $m\ll n^2$ 时，邻接矩阵中将包含大量无效元素，从而造成明显的空间浪费.

因此，在实际工程实现中，更常采用适用于稀疏图的邻接表及其压缩存储形式.本文主要涉及两类常见的稀疏存储结构，即 COO（Coordinate）格式与 CSR（Compressed Sparse Row）格式.

其中，COO 格式使用三个长度为 $m$ 的一维数组分别记录每条边的起点索引、终点索引以及对应权值.该格式结构直观、实现简单，适合用于图结构的构造与调试.

与之相比，CSR 格式在 COO 的基础上进一步压缩了行索引信息.它通常使用 \texttt{row\_ptr}、\texttt{col\_indices} 和 \texttt{value} 三个数组来描述图结构，其中 \texttt{row\_ptr[i]} 与 \texttt{row\_ptr[i+1]} 共同确定了第 $i$ 个顶点所有出边在 \texttt{col\_indices} 与 \texttt{value} 中对应的存储区间.相比 COO，CSR 在按顶点批量访问邻接边时更加高效，同时能够在保证结构信息完整的前提下进一步节省空间.图~\ref{fig:csr_storage} 展示了邻接矩阵 $C$ 与 CSR 存储格式之间的对应关系.可以看出，CSR 并未改变图原有的拓扑结构，而只是对其进行了更紧凑的线性化表达.由于本文后续涉及线图构建、邻接展开以及 GPU 并行访存等操作，CSR 形式也更符合后续算法实现的需求.

\begin{figure}[H]
  \centering
  \begin{tikzpicture}[
    >=Latex,
    font=\small,
    box/.style={draw,minimum height=6mm,minimum width=7mm,inner sep=1.5pt},
    label/.style={anchor=east},
    arr/.style={->,line width=0.45pt}
  ]
    % Left: adjacency matrix C
    \matrix (C) [matrix of math nodes, left delimiter={(}, right delimiter={)},
      nodes={inner sep=1.2pt}, row sep=2.5pt, column sep=4.5pt] {
      0 & 2 & 0 & 7 \\
      0 & 0 & 3 & 0 \\
      4 & 0 & 0 & 1 \\
      0 & 5 & 0 & 0 \\
    };
    \node[left=3mm of C.west] {$C=$};

    % Anchor for CSR block (shift right to avoid overlaps)
    \coordinate (csr) at ([xshift=50mm]C.east);
    \draw[arr] ([xshift=4mm]C.east) -- ([xshift=26mm]C.east);

    % Right: CSR arrays
    \node[label] (rpl) at ([yshift=9mm]csr) {\texttt{row\_ptr}:};
    \node[box] (rp0) at ([xshift=14mm,yshift=9mm]csr) {0};
    \node[box] (rp1) [right=0pt of rp0] {2};
    \node[box] (rp2) [right=0pt of rp1] {3};
    \node[box] (rp3) [right=0pt of rp2] {5};
    \node[box] (rp4) [right=0pt of rp3] {6};

    \node[label] (cil) at ([yshift=1mm]csr) {Col\_indices:};
    \node[box] (ci0) at ([xshift=14mm,yshift=1mm]csr) {1};
    \node[box] (ci1) [right=0pt of ci0] {3};
    \node[box] (ci2) [right=0pt of ci1] {2};
    \node[box] (ci3) [right=0pt of ci2] {0};
    \node[box] (ci4) [right=0pt of ci3] {3};
    \node[box] (ci5) [right=0pt of ci4] {1};

    \node[label] (val) at ([yshift=-7mm]csr) {Value:};
    \node[box] (v0) at ([xshift=14mm,yshift=-7mm]csr) {2};
    \node[box] (v1) [right=0pt of v0] {7};
    \node[box] (v2) [right=0pt of v1] {3};
    \node[box] (v3) [right=0pt of v2] {4};
    \node[box] (v4) [right=0pt of v3] {1};
    \node[box] (v5) [right=0pt of v4] {5};
  \end{tikzpicture}
  \caption{CSR 格式存储示意图}
  \label{fig:csr_storage}
\end{figure}

\section{最短路径问题}


\subsection{单源最短路径问题的定义}

\begin{definition}
在图 $G(V, E, W)$ 中，若边集合 $(v_i, v_{i+1}), (v_{i+1}, v_{i+2}), \ldots, (v_{p-1}, v_j)$ 存在，则序列 $p = \langle v_i, v_{i+1}, \ldots, v_j \rangle$ 称为从点 $v_i$ 到点 $v_j$ 的一条路径.为方便表述复杂网络中的边，图的单条边也常由希腊字母 $\alpha, \beta$ 等直接表示.路径 $p$ 的距离（或代价值）定义为
\[
f(p) = \sum_{k=i}^{j-1} w(v_k, v_{k+1}).
\]
\end{definition}

\begin{definition}
最短路径是指从起点 $s$ 到点 $v_j$ 的所有路径中，使得路径长度（即所有边权重之和）最小的路径.若存在多条最短路径，则任选其一作为最短路径.最短路径长度为
\[
d(v_i, v_j) = \min(f(p)),
\]
其中 $p$ 为所有从 $v_i$ 到 $v_j$ 的路径集合中的任意路径.
\end{definition}

\begin{definition}
设 $u$ 和 $v$ 是图 $G(V, E, W)$ 中两个顶点，$w(u,v)$ 表示顶点 $u$ 到点 $v$ 的边的权重，$d(u)$ 和 $d(v)$ 分别表示源点 $s$ 到顶点 $u$ 和顶点 $v$ 的最短距离.如果 $d(u) + w(u,v) < d(v)$，则更新 $d(v)$ 为 $d(u) + w(u,v)$.这一过程就是松弛操作.
\end{definition}

单源最短路径问题即求源点到所有顶点的最短路径，核心操作为反复松弛.

\subsection{经典最短路径问题算法}

随着对最短路径问题研究的不断深入，相关领域提出了多种经典求解方法.其中，Dijkstra 算法作为早期最具代表性的经典算法之一，在边权非负的单源最短路径问题中具有重要地位；而 $\Delta$-Stepping 算法则在传统最短路径算法基础上进一步发展而来，兼具较好的求解效率与并行实现潜力.因此，本节将围绕 Dijkstra 算法与 $\Delta$-Stepping 算法展开介绍，并分析其基本思想、实现特点及适用条件.

\subsubsection{Dijkstra 算法}

Dijkstra算法由荷兰计算机科学家E.W. Dijkstra于1956年提出，是解决单源最短路径问题最常用的贪心算法.该算法适用于所有边权重均为非负的有向图或无向图.


算法思想：Dijkstra算法采用贪心策略，每次从未访问的顶点中选择距离源点最近的顶点，将其加入已访问集合，并据此更新其他顶点的最短距离.通过对所有顶点的这样处理，最终可得到从源点到所有其他顶点的最短路径.


时间复杂度：采用二叉堆实现的Dijkstra算法时间复杂度为 $O((|V| + |E|) \log |V|)$；采用线性查找实现的Dijkstra算法时间复杂度为 $O(|V|^2)$.


算法伪代码：

\begin{algorithm}[H]
  \SetAlgoLined
  \caption{Dijkstra Algorithm}
  \label{algo:dijkstra}
  
  Input: Graph $G = (V, E, W)$, source vertex $s$\\
  Output: Shortest distances $d(v)$ for all $v \in V$\\
  
  Initialization:\\
  \For{each $v \in V$}{
    $d(v) \leftarrow \infty$\
  }
  $d(s) \leftarrow 0$\
  $T \leftarrow V/s$\
  
  \While{$T$ is not empty}{
    Min: $u \leftarrow \{v : v \in T \text{ and } \forall c \in T, d(v) \leq d(c)\}$\
    
    \If{$d(u) = \infty$}{
      break\
    }
    
    Remove $u$ from $T$\
    
    Relax:\
    \For{each $(u, v) \in E$}{
      \If{$d(u) + w(u, v) < d(v)$}{
        $d(v) \leftarrow d(u) + w(u, v)$\
      }
    }
  }
  
  \Return{$d$}\
\end{algorithm}

在上述伪代码中，$d(v)$ 表示源点 $s$ 到顶点 $v$ 的最短距离，$T$ 表示未处理顶点的集合，$w(u, v)$ 表示边 $(u, v)$ 的权重.

算法步骤说明：
\begin{enumerate}
  \item 初始化：将所有顶点的距离初始化为无穷大，源点距离为0，未处理顶点集合 $T$ 为除源点外的所有顶点.
  \item 主循环：当存在未处理的顶点时，重复以下步骤：
  \begin{enumerate}
    \item 从未处理顶点集合 $T$ 中选择距离最小的顶点 $u$.
    \item 若 $d(u) = \infty$，说明剩余顶点不可达，退出循环.
    \item 将顶点 $u$ 从集合 $T$ 中移除.
    \item 对从 $u$ 出发的所有边执行松弛操作，若发现更短路径则更新距离.
  \end{enumerate}
  \item 输出：返回所有顶点的最短距离数组 $d$.
\end{enumerate}

优缺点分析：
\begin{itemize}
  \item 优点：计算效率高，时间复杂度为 $O((|V|+|E|)\log|V|)$（使用优先队列实现）；适用于边权非负的单源最短路径问题.
  \item 缺点：无法处理带有负权边的图；对于稠密图，优先队列的操作可能导致性能下降.
  \item 与非加性问题的联系：Dijkstra算法依赖于边权的加性性质，无法直接应用于非加性问题，需要对问题进行转化或设计新的算法.
\end{itemize}

\subsubsection{\texorpdfstring{$\Delta$}{Delta}-Stepping 算法}

为了解决 Dijkstra 算法中严格的优先队列带来的难以并行化问题，Meyer 等人于 2003 年提出了 $\Delta$-Stepping 算法.该算法通过引入“桶（Bucket）”数据结构，在 Dijkstra 算法与 Bellman-Ford 算法之间取得了良好的平衡，从而实现了对图数据结构的有效并行处理.

\textbf{算法思想}：$\Delta$-Stepping 算法的核心在于将优先队列离散化为一系列连续的桶数组 $B$.设定一个步长参量 $\Delta > 0$，每个桶 $B_i$ 用于存放当前暂定距离 $d(v)$ 落在区间 $[i\Delta, (i+1)\Delta)$ 范围内的顶点.在每一轮迭代中，算法并发地从当前非空的编号最小的桶中取出所有顶点，并使用它们的临接“轻边”（权重 $\le \Delta$ 的边）去更新其后继顶点的距离.如果更新后的顶点距离落入了更小的桶区间，则将其移入对应的桶中.当当前桶内的顶点经过轻边完全松弛稳定后，再利用与之相连的“重边”（权重 $> \Delta$ 的边）进行一轮跨桶范围的并行扩展松弛.

\textbf{时间复杂度}：对于带最大度数限制的随机图拓扑结构，利用并行模型实现的 $\Delta$-Stepping 算法其平均时间复杂度可趋近于 $O(|V| + |E| + \frac{L}{\Delta})$，其中 $L$ 为最长最短路径距离.

\textbf{在非加性并行寻路中的应用机理}：在应对如线图高阶网格这种极度稠密网络且由于需要处理千万级节点状态的背景下，$\Delta$-Stepping 算法有效地避免了并行线程对于单一队列的锁争用（Lock Contention）.通过控制共享变量的并发原子访问（如使用 \texttt{atomicMin}），本算法可以在现代 GPU 计算管线上直接映射为无强制排序的桶内波前探索操作，这也是其在海量数据非加性路由重构系统内进行寻优加速的直接理论基础.

\section{非加性最短路径问题}

非加性最短路径问题放宽了加性假设，使边权可能随上下文、时间或约束变化，导致经典算法失效.本节说明定义、失效原因、现实普适性，并据此给出分类与求解思路.

\subsection{定义与性质}

最短路径的核心是路径成本函数.若路径成本可写成各边固定权重的算术和，则属于加性最短路径；否则即为非加性最短路径.

\subsubsection{加性最短路径问题的定义}

\begin{definition}[加性最短路径问题]
设路径 $p = \langle v_0, v_1, v_2, \ldots, v_k \rangle$ 为从源点 $v_0 = s$ 到终点 $v_k = t$ 的路径，若路径的总成本满足
\[
    f(p) = \sum_{i=0}^{k-1} w(v_i, v_{i+1}),
\]
即“等于路径所经过各边权重的算术和”，且“各边的权重为恒定的、与路径中其他边无关的固定值”，则称该最短路径问题为加性最短路径问题（Additive Shortest Path Problem, ASPP）.
\end{definition}

加性问题的核心特征包括：

（1）边权的独立性和恒定性：在加性最短路径问题中，每条边 $(u, v)$ 的权重 $w(u, v)$ 是预定义的、恒定的常数，与该边在路径中的位置、路径中的其他边、或任何其他上下文无关.这是加性问题最本质的假设，保证了成本的可预测性与确定性.

（2）无后效性（Memoryless Property）：子路径的成本独立于路径的后续扩展，即对任意路径 $p_1 = \langle s, \ldots, u \rangle$ 和 $p_2 = \langle u, \ldots, t \rangle$，有
\[
    f(p_1 \oplus p_2) = f(p_1) + f(p_2),
\]
其中 $\oplus$ 表示路径拼接.历史信息（前面的路径）对未来成本无影响.

（3）最优子结构（Optimal Substructure）：最优路径必然由各段局部最优子路径构成.若 $p^* = \langle s, \ldots, u, \ldots, t \rangle$ 是从 $s$ 到 $t$ 的最优路径，则 $p^*$ 中从 $s$ 到 $u$ 的子路径必为 $s$ 到 $u$ 的最优路径.这直接导出递推关系
\[
    d(t) = \min\{d(u) + w(u, t) : (u, t) \in E\},
\]
这是 Dijkstra、Bellman-Ford 等经典贪心与动态规划算法高效求解的理论根基.

\subsubsection{非加性最短路径问题的定义}

\begin{definition}[非加性最短路径问题]
若路径总成本 $f(p)$ 不能表示为
\[
  f(p) = \sum_{i=0}^{k-1} w(v_i, v_{i+1}),
\]
或边权不再恒定、独立（随时间、位置或历史决策变化），则该最短路径问题称为非加性最短路径问题（Non-additive Shortest Path Problem, NASPP）.
\end{definition}

非加性问题破坏了加性假设的至少一条，带来三类失效：

边权动态性与耦合性：权重依赖时间、路径位置或先前选择，如时变网络 $w(u,v,t)$、容量受限网络的拥塞反馈、多目标场景的向量权重 $\mathbf{w}(u,v)$.

无后效性失效：历史决策影响后续成本，例如车辆疲劳模型、带宽占用导致后续代价（如转弯代价等）依赖当前及前驱路径 $\alpha, \beta$.

最优子结构破坏：全局最优不能由子问题最优递推，如极值型目标 $f(p)=\max w(v_i,v_{i+1})$、多目标 Pareto 优化、带全局约束的可行性组合.

\subsection{问题分类}

非加性最短路径问题在实际应用中具有广泛来源，其表现形式也较为多样.与传统加性最短路径问题相比，这类问题的共同特点在于：路径总代价不再单纯由各边固定权重的线性累加所决定，而往往会受到约束条件、环境状态、路径历史或多属性耦合关系的影响.为了便于后文针对不同问题特征选择相应的建模与求解方法，本节对较为常见的几类非加性最短路径问题进行归纳和说明.需要指出的是，本文后续重点讨论的是路径依赖型非加性最短路径问题，因此其余几类问题主要作为背景性分类进行概述，不作深入展开.

\subsubsection{瓶颈路径问题}

瓶颈路径问题（Bottleneck Path Problem）并不以路径总代价最小作为优化目标，而是关注路径上的“最薄弱环节”或“关键限制因素”.其典型形式包括最小化路径上的最大边权，即 min-max 问题；以及最大化路径上的最小边权，即 max-min 问题，后者也常被称为最宽路径问题.这类模型通常用于描述带宽路由、运输载重限制、桥梁限高、最大风险暴露控制等应用场景.在这些问题中，决定路径优劣的往往不是各段代价的总和，而是其中最不利的一段边或最受限制的一类资源，因此其目标函数本质上不满足经典最短路径中的线性可加结构.尽管如此，由于瓶颈目标具有较强的结构特征，在某些特定情形下仍可借助生成树、阈值判定或改造后的贪心方法实现多项式时间求解.从分类角度看，这类问题说明：一旦路径评价标准由“总和最优”转变为“极值最优”，最短路径问题便已偏离传统加性框架\cite{pollack1960maxcap,edmonds1972theoretical,camerini1984random,gj1979np}.

\subsubsection{资源约束型最短路径问题}

资源约束型最短路径问题（Resource Constrained Shortest Path Problem, RCSP）是在路径成本最优化的基础上，进一步引入时间窗、电量、容量、预算、服务时限等附加资源约束的一类问题.在这类问题中，路径是否可行不仅取决于当前所处节点，还取决于此前路径上已经消耗了多少资源，因此路径的可行性与历史累计状态之间存在紧密联系.也就是说，即便两条路径到达同一节点，只要它们在资源使用情况上不同，后续可扩展性也可能完全不同.这种特征决定了问题求解过程中通常需要同时维护“成本 + 资源状态”的复合信息，而不能仅用单一距离值描述节点状态.随着资源维度的增加，状态空间往往迅速膨胀，导致问题复杂度显著提升，因此该类问题在一般形式下通常属于 NP-hard.实际求解中，常见方法包括标签算法、动态规划剪枝、列生成以及整数规划等.资源约束型问题表明，非加性的一个重要来源，是“路径的历史消耗”对后续决策空间的持续影响\cite{pugliese2013survey,irnich2005spprc,barnhart1998branch,bode2012cutfirst}.

\subsubsection{多属性非线性组合最短路径问题}

多属性非线性组合最短路径问题是指路径优劣不再由单一指标决定，而是需要同时综合时间、费用、风险、能耗、舒适性等多个属性进行评价.在许多现实应用中，这些属性之间并不是简单的线性加权关系，而常常通过某种非线性效用函数、风险函数或偏好函数进行组合.例如，用户对“更快但更贵”与“更便宜但更慢”的选择偏好，往往不能仅通过固定线性权重准确刻画；又如，当风险随累计暴露时间非线性增加时，路径总代价也不再是边权的简单叠加.在这种情况下，问题的目标函数通常不具备经典最短路径所要求的可分解性和最优子结构，因而难以直接应用 Dijkstra 等传统算法.该类问题通常需要结合多目标优化、鲁棒优化、效用建模或 Pareto 最优分析等方法进行处理.从本质上看，这类问题强调的是“多维属性之间的非线性交互”如何打破路径代价的线性累加机制\cite{ehrgott2005multicriteria,deb2002nsga2,kim2006adaptive,bertsimas2004price}.

\subsubsection{时间依赖型最短路径问题}

时间依赖型最短路径问题（Time-Dependent Shortest Path Problem, TDSP）是指边权会随出发时刻或到达时刻变化的一类路径优化问题.典型例子包括城市交通网络中的拥堵变化、公共交通系统中的发车时刻表、机场航班衔接网络以及物流调度中的动态通行时间等.在这类问题中，同一条边在不同时间通过时可能具有完全不同的代价，因此边权不再是静态常数，而是一个与时间状态相关的函数.这样一来，路径总代价不仅依赖于所选边集本身，还依赖于路径经过这些边时所对应的具体时刻.由于“到达当前边的时间”本身又由此前路径决定，因此路径成本会表现出明显的动态耦合特征.时间依赖型问题的难点在于，路径的局部选择会改变后续边的实际代价，从而使得传统最短路径中的静态边权假设失效.对于此类问题，研究中常借助时间扩展图、时间标签更新以及满足 FIFO 条件下的动态最短路径算法进行建模和求解\cite{cooke1966shortest,orda1990shortest,dean2005fifo,chabini1998discrete}.

\subsubsection{随机性与相关性最短路径问题}

在一些更复杂的实际网络中，边权并不是确定的固定值，而可能表现为随机变量，并且不同边之间还可能存在相关结构.由此产生的随机性与相关性最短路径问题（Stochastic and Correlated Shortest Path Problem, SCSP）通常用于描述具有不确定性的交通网络、通信网络或供应链系统.例如，道路旅行时间可能受到天气、事故和整体交通流状态影响，从而在不同路段之间呈现相关波动；又如，某些网络中的链路失效概率也可能并非相互独立.在这类问题中，路径优化目标可能是最小期望时间、最大到达可靠性、最小方差风险，或者在给定置信水平下最小化某种风险度量.由于需要同时处理概率分布、随机波动以及边间相关性，路径代价的表达方式显然已不满足简单的确定性加法结构.相应地，这类问题在建模和求解上通常需要引入随机规划、鲁棒优化、采样近似或风险约束等方法.该类问题说明，非加性不仅可能来源于路径历史，也可能来源于环境不确定性及边权之间的统计依赖关系\cite{frank1969probabilistic,nikolova2006approx,birge2011stochastic,bertsimas2011robust}.

\subsubsection{路径依赖型最短路径问题}

路径依赖型最短路径问题是本文重点关注的一类非加性最短路径问题.其核心特征在于，当前一步的代价不仅由当前边本身决定，还会受到此前路径状态的影响，例如前一条边、已访问节点序列、交通转向状态、模式切换状态，或者累计资源消耗等.换言之，在这类问题中，系统的真实状态已经不再只是“当前所在节点”，而应当表示为“当前节点 + 历史状态”的组合.这意味着，即便两条路径到达同一个节点，只要它们的历史不同，后续扩展所产生的代价和可行性也可能不同.正是这种对历史信息的依赖，使得问题的状态空间迅速膨胀，也使得传统基于节点最短距离递推的经典最短路框架难以直接适用.

在路径依赖型问题中，一类最常见、也最具有代表性的情形，是当前转移代价依赖于前一条边，即代价定义在“相邻边对”之上.例如，在道路网络中，车辆从一条边转入下一条边时，可能产生转向惩罚、禁转约束、等待成本或模式切换成本；在这类情况下，代价已经不再属于单条边，而是属于边与边之间的过渡关系.对于这种局部历史依赖问题，线图转换提供了一种较为自然且结构清晰的建模方式：将原图中的边映射为线图中的节点，再将原图中“相邻两条边可以连续连接”的关系映射为线图中的边.经过这一变换，原本依附于相邻边对的代价，可以转化为线图上的静态边权，从而把原问题重新组织为适合标准最短路径框架处理的形式.也正因为如此，本文后续将重点围绕这一路径依赖型问题展开讨论，并进一步分析线图方法在该类问题中的适配性与并行实现潜力\cite{gj1979np,thomas2007timepath,irnich2005spprc,desaulniers2014vrptw}.

\subsection{路径依赖型最短路径问题的经典算法}

在非加性最短路径问题的六大分类中，路径依赖型最短路径问题具有显著的代表性和研究价值.其核心特征是：当前边或转移的成本不仅取决于边本身，还取决于之前经过的路径历史（如已访问节点序列、累计资源消耗、转向惩罚等）.由于路径依赖性的引入，问题状态空间从单一节点扩展为“节点 + 历史路径/状态”的组合，导致状态空间指数级增长，最优子结构性质失效，问题复杂度通常为 NP-hard\cite{gj1979np}.鉴于本文后续研究将聚焦于此类问题，本节重点介绍两类经典求解算法，并分析其局限性，为第三章线图转换方法的引入奠定基础.

\subsubsection{标签算法（Labeling Algorithm）}

标签算法是求解路径依赖型最短路径问题的经典精确方法，其基本思想是为每个节点维护多个标签来表示"到达该节点的不同历史状态"\cite{irnich2005spprc}.为便于描述，设原图 $G=(V,E)$，源点为 $s$，目标点为 $t$.一个标签可抽象为
\[
  \ell = (v,\; g,\; \mathbf{r},\; h),
\]
其中 $v\in V$ 表示当前节点，$g$ 表示累计成本，$\mathbf{r}$ 表示资源消耗（可为多维），$h$ 表示关键历史信息（例如已访问节点集合、上一步边/转向状态等）.标签算法通过扩展标签生成新标签，并利用支配规则进行剪枝：若在同一节点 $v$ 上，标签 $\ell_A$ 在成本与资源等维度上均不劣于 $\ell_B$，且其历史信息不比 $\ell_B$ 更“受限”（从而任意后续扩展都不会更差），则称 $\ell_A$ 支配 $\ell_B$，可安全删除 $\ell_B$.标签算法的伪代码如算法\ref{algo:pdsp_labeling}所示.

\begin{algorithm}[H]
  \SetAlgoLined
  \caption{Labeling Algorithm for Path-Dependent Shortest Path}
  \label{algo:pdsp_labeling}

  Input: Graph $G=(V,E)$, source $s$, target $t$; transition cost $c(u,v,h)$; state update $h'\!=\!\tau(h,u,v)$; feasibility check $\textsc{Feasible}(\mathbf{r}')$\\
  Output: Best label reaching $t$ (and its path)\\

  Initialization:\\
  For each $v\in V$, set label set $\mathcal{L}(v) \leftarrow \emptyset$\\
  Create initial label $\ell_0=(s,0,\mathbf{0},h_0)$ and insert into $\mathcal{L}(s)$\\
  Initialize a queue (or priority queue) $Q\leftarrow\{\ell_0\}$\\

  \While{$Q$ is not empty}{
    Pop a label $\ell=(u,g,\mathbf{r},h)$ from $Q$\\
    \If{$\ell$ is dominated by some label in $\mathcal{L}(u)$}{
      continue\\
    }
    \For{each outgoing edge $(u,v)\in E$}{
      $h' \leftarrow \tau(h,u,v)$\\
      $\mathbf{r}' \leftarrow \textsc{UpdateRes}(\mathbf{r},u,v,h)$\\
      \If{$\textsc{Feasible}(\mathbf{r}')$}{
        $g' \leftarrow g + c(u,v,h)$\\
        Create new label $\ell'=(v,g',\mathbf{r}',h')$\\
        \If{$\ell'$ is not dominated by labels in $\mathcal{L}(v)$}{
          Remove from $\mathcal{L}(v)$ all labels dominated by $\ell'$\\
          Insert $\ell'$ into $\mathcal{L}(v)$ and push $\ell'$ into $Q$\\
        }
      }
    }
  }

  \Return{the best (minimum-$g$) label in $\mathcal{L}(t)$}\\
\end{algorithm}

伪代码说明：
\begin{itemize}
  \item 标签的含义：同一节点 $v$ 上的不同标签对应不同的“到达方式”，它们携带了导致后续代价差异的历史信息 $h$ 与资源状态 $\mathbf{r}$.
  \item 扩展与更新：沿边 $(u,v)$ 扩展时，用 $\tau(\cdot)$ 更新历史状态（例如更新“上一条边/转向方向”或把访问节点加入集合），并用 $\textsc{UpdateRes}$ 更新资源消耗；同时用 $c(u,v,h)$ 计算依赖历史的边代价.
  \item 支配剪枝：在 $\mathcal{L}(v)$ 内做支配判断，保留“潜在更优且更可扩展”的标签，删除明显不可能产生最优解的标签，以减少搜索空间.
\end{itemize}

方法局限性：尽管标签算法可保证在给定支配规则与状态刻画下找到全局最优解，但标签数量往往随路径长度呈指数增长，导致内存消耗巨大；支配判断本身在多维资源与复杂历史编码下代价较高；并且标签扩展具有强数据相关性，整体过程偏串行，难以直接高效并行化，因而在大规模图或实时场景下计算效率受限.

\subsubsection{列生成与分支定价方法（Column Generation / Branch-and-Price）}

除直接在节点上维护多标签外，路径依赖型最短路径问题在运筹优化中还常被放入“列生成（column generation）/分支定价（branch-and-price）”框架中求解.其核心思想是：不显式枚举所有可行路径，而是把“从 $s$ 到 $t$ 的一条可行路径”视为一个变量（列），在主问题中选择若干列使目标最小；随后通过定价子问题按需生成能够改进目标的路径列.

形式上，可将每条可行路径 $p\in\mathcal{P}$ 赋予成本 $C(p)$（其中 $C(\cdot)$ 可依赖历史/资源/状态演化），令二元变量 $x_p\in\{0,1\}$ 表示是否选择路径 $p$.在只选一条 $s$--$t$ 路径的基本情形下，主问题可写为
\[
  \min\ \sum_{p\in\mathcal{P}} C(p)\,x_p,\qquad
  \text{s.t.}\ \sum_{p\in\mathcal{P}} x_p = 1,\ x_p\ge 0\ (\text{或 }x_p\in\{0,1\}).
\]
由于 $|\mathcal{P}|$ 往往极大，列生成从一个小的路径集合 $\mathcal{P}_0$ 开始求解受限主问题（RMP），并根据对偶信息构造“约化成本（reduced cost）”来寻找新的改进列；若考虑整数性，则在分支过程中反复进行列生成，即形成分支定价框架.

\begin{algorithm}[H]
  \SetAlgoLined
  \caption{Column Generation Framework for Path-Dependent Shortest Path}
  \label{algo:pdsp_column_generation}

  Input: Feasible path set $\mathcal{P}$ (implicit), path cost $C(p)$; initial columns $\mathcal{P}_0$\\
  Output: Best path (integral if solved with branching)\\

  Initialization:\\
  Set working columns $\mathcal{P}_R \leftarrow \mathcal{P}_0$\\

  \Repeat{no improving column found}{
    Solve Restricted Master Problem (RMP) over $\mathcal{P}_R$ and get duals $\pi$\\
    Pricing: find a feasible path $p^*\in\mathcal{P}$ minimizing reduced cost $\bar C(p)=C(p)-\pi$\\
    \If{$\bar C(p^*) < 0$}{
      Add $p^*$ into $\mathcal{P}_R$\\
    }
  }

  \If{integrality is required}{
    Apply branching rules and repeat column generation in nodes (branch-and-price)\\
  }

  \Return{best path implied by final solution}\\
\end{algorithm}

伪代码说明：该框架将“求解路径依赖最短路径”转化为“迭代求解一个小规模主问题 + 反复解定价子问题”.其中定价子问题本质上是在路径依赖约束/状态演化下寻找约化成本最小的可行路径，常见做法是复用标签扩展与支配剪枝等机制在隐式状态空间中搜索，从而避免枚举全部路径.

方法局限性：列生成/分支定价可在大规模组合优化中获得较强的求解能力，但其计算代价通常集中在反复求解定价子问题与分支树扩展上；当路径依赖导致可行路径结构高度复杂时，定价问题本身可能很难、且需要大量标签与支配判断；此外主问题与子问题之间存在强耦合迭代，整体流程同样具有明显串行瓶颈.

\subsubsection{算法局限性与线图转换的动机}

总体而言，无论是直接在图上进行标签扩展与支配剪枝，还是采用列生成/分支定价这类“主问题--定价子问题”的迭代框架，核心困难都绕不开两点：其一是路径依赖带来的隐式状态空间膨胀，使得标签数、可行路径数或搜索分支数迅速增长，从而带来巨大的时间与内存压力；其二是扩展、支配判断与迭代求解具有强数据相关性与顺序依赖，导致并行化与工程落地难度较高.

针对上述局限，本文引入线图（Line Graph）转换作为后续的求解思路：线图将原图的边转化为节点，并将“相邻边之间的依赖关系”显式编码为线图节点间的连接关系，从而把部分路径依赖信息结构化地映射到图结构中.相较于在隐式状态空间中维护大量标签或反复求解定价问题，线图转换在规模增长的可控性、结构规则性以及并行实现方面更具潜力.第三章将系统阐述线图转换的理论基础、算法设计与并行优化策略.

\section{线图理论基础}

在路径依赖型最短路径问题（Path-dependent Shortest Path Problem）中，路径的整体代价不仅取决于单条边的独立权重，还与相邻边之间的组合与连接状态密切相关.传统的图拓扑结构以节点为核心，难以直接刻画边与边之间的动态耦合关系.为了将这种隐藏的“边关系”显式化，引入结构转换机制成为一种有效的解决途径.线图（Line Graph）作为一种经典的图转换方法，通过实现图元素角色的对偶转换，已在图论及复杂网络优化领域得到广泛应用\cite{whitney1932congruent,harary1960some,van1965interchange}.本节将系统梳理线图的基本定义、构造范式及其拓扑性质，为后续求解算法的设计提供坚实的理论支撑.

\subsection{线图的定义}

设原图为 $G=(V,E)$，其对应的线图记为 $L(G)=(V',E')$\cite{whitney1932congruent}.线图的基本思想可概括为“化边为点”：原图中的每一条边 $e\in E$，在转换后的线图中对应为一个顶点；而原图中两条能够首尾相接的边，则在线图中对应为一条连接边.

对于有向图而言，上述规则可以进一步具体化\cite{harary1960some}.设原图中存在两条有向边 $e_1=(u,v)$ 与 $e_2=(v,w)$，若 $e_1$ 的终点恰好是 $e_2$ 的起点，则说明这两条边在原图中可以连续经过.于是，在线图 $L(G)$ 中，便在表示 $e_1$ 与 $e_2$ 的两个顶点之间建立一条有向边 $(e_1,e_2)$.

因此，线图的本质并不是简单复制原图，而是将原图中的“边关系”重新组织为线图中的“点关系”.这一转换使得原本需要借助路径上下文才能识别的边间衔接关系，在线图中被直接编码为显式的邻接结构，为后续刻画边间依赖提供了更直接的表示方式.

\subsection{线图的构造过程}

对于任意有向图 $G=(V,E)$，其线图 $L(G)$ 的构造过程通常可概括为以下几个步骤.

首先，遍历原图中的每一条有向边 $e=(u,v)$，并在线图中为其创建一个对应顶点.换言之，线图的顶点集合由原图的边集合诱导而来.

其次，对原图中任意一对边 $e_1=(u,v)$ 与 $e_2=(v,w)$ 进行检查.若二者满足前一条边的终点等于后一条边的起点，则说明它们在原图中能够构成合法的连续转移，此时在线图中加入一条从 $e_1$ 指向 $e_2$ 的有向边.

最后，对所有满足连接条件的边对重复上述操作，即可得到完整的线图拓扑结构.图~\ref{fig:line_graph_example} 展示了这一转换过程的示意：原图 $G^0$ 的边 $(1,2)$ 与 $(2,4)$ 在线图 $G^1=L(G^0)$ 中分别对应为点 $12$ 与 $24$，并由一条有向边连接.通过该变换，原图中“边与边之间的衔接”被展开为线图中“点与点之间的连接”，从而为后续处理边间依赖提供了更直接的表示方式.

\begin{figure}[H]
  \centering
  \begin{tikzpicture}[
    >=Latex,
    main node/.style={circle,draw,minimum size=6.5mm,inner sep=0pt},
    edge label/.style={draw=none,fill=white,inner sep=1pt,font=\footnotesize},
    every edge/.style={draw,->,line width=0.6pt}
  ]
    % Original Graph G
    \begin{scope}[xshift=0cm, yshift=0cm]
      \node[main node] (U)  at (0, 0) {u};
      \node[main node] (X)  at (0, -2.0) {x};
      \node[main node] (V)  at (2.5, -1.0) {v};
      \node[main node] (W1) at (5.0, 0.5) {w$_1$};
      \node[main node] (W2) at (5.0, -2.5) {w$_2$};
      
      \path (U) edge node[edge label, above=2pt, sloped] {$e_1$} (V)
            (X) edge node[edge label, below=2pt, sloped] {$e_4$} (V)
            (V) edge node[edge label, above=2pt, sloped] {$e_2$} (W1)
            (V) edge node[edge label, below=2pt, sloped] {$e_3$} (W2);
            
      \node[draw=none] at (2.5, -3.5) {\textbf{原图:} $G$};
    \end{scope}
    
    % Arrow indicating transformation aligned vertically with the center (-1.0)
    \node[draw=none] (arrow) at (6.6, -1.0) {\huge $\Rightarrow$};
    
    % Line Graph L(G)
    \begin{scope}[xshift=8.5cm, yshift=0cm]
      \node[main node] (E1) at (0, 0) {$e_1$};
      \node[main node] (E4) at (0, -2.0) {$e_4$};
      \node[main node] (E2) at (3.0, 0.5) {$e_2$};
      \node[main node] (E3) at (3.0, -2.5) {$e_3$};
      
      % Shift positions of diagonal labels so they don't visually intersect
      \path (E1) edge node[edge label, above=1pt, sloped] {$(e_1, e_2)$} (E2)
            (E1) edge node[edge label, pos=0.74, above=1pt, sloped] {$(e_1, e_3)$} (E3)
            (E4) edge node[edge label, pos=0.26, above=1pt, sloped] {$(e_4, e_2)$} (E2)
            (E4) edge node[edge label, below=1pt, sloped] {$(e_4, e_3)$} (E3);
      
      \node[draw=none] at (1.5, -3.5) {\textbf{线图:} $L(G)$};
    \end{scope}

  \end{tikzpicture}
  \caption{有向图网络拓扑映射至线图的转换展开结构}
  \label{fig:line_graph_example}
\end{figure}


\subsection{线图的拓扑性质}

线图在拓扑结构层面呈现出若干关键性质，这些性质是后续算法设计的核心依据：

（1）拓扑规模扩张性

线图的节点规模严格等于原图的边规模，即：
$$ |V(L(G))| = |E(G)| $$ 
由于实际网络（如交通网、通信网）通常满足 $|E| \gg |V|$，线图变换会导致图规模的非线性扩张.这一特性提示我们在算法设计时，必须考量由于状态空间膨胀带来的存储与计算开销问题.

（2）邻接语义的重定义

在线图空间中，“邻接关系”不再表征地理或物理节点的直接连通，而是刻画原图中不同路段之间的“时空衔接连贯性”.这使得线图具备了捕获原图高阶路径结构特征的能力.

（3）路径结构的同构映射

原图空间中的任意一条可行路径：
$$ P_G = v_1 \rightarrow v_2 \rightarrow \cdots \rightarrow v_k $$ 
在线图空间中存在唯一对应的一条路径：
$$ P_{L(G)} = (v_1,v_2) \rightarrow (v_2,v_3) \rightarrow \cdots \rightarrow (v_{k-1},v_k) $$ 
这种严格的同构映射不仅保证了路径连通性的无损传递，也确保了在双域间进行最优性检索的等价性.

（4）依赖信息的显式表征优势

由于线图的顶点直接关联原图的边，研究者可在线图顶点或连边上灵活挂载多维属性矩阵.例如，可将转弯延误、路口惩罚或复杂的边际交互成本直接编码为线图的边权.这一性质从根本上消除了经典图结构在处理“非马尔可夫型”路径成本时的表达瓶颈.

\subsection{线图在最短路径求解中的应用机理}

在经典最短路径范式中，路径评价体系通常遵循严格的线性可加性假设.然而，在面对包含复杂转弯限制或状态相依成本的非加性网络时，贝尔曼最优性原理（Bellman's Principle of Optimality）往往失效，导致 Dijkstra 等传统算法难以直接适配.

引入线图理论可实现问题域的降维重构.具体映射机理如下：

\begin{itemize}
  \item 状态转移：将原图难以处理的“边间耦合约束”转化为线图中常规的“节点与边权属性”；
  \item 成本线性化：将非加性的路径依赖成本吸收进线图的拓扑结构中，使得转换后的线图重新满足权值累加原则.
\end{itemize}

经过此类同构变换，复杂的路径依赖型最短路径问题被优雅地还原为线图空间下的标准最短路径问题.此时，不仅传统的高效路由算法能够重新发挥效用，且线图本身高度规则的矩阵特征也为引入现代并行计算框架（如 GPU 加速）创造了有利条件.

值得注意的是，若非加性成本不仅仅取决于前一条边，而是依赖于前 $k$ 条边的路径历史，则此类问题可以通过构造高阶线图（Higher-order Line Graph，记为 $L^{[k]}(G)$）来解决.高阶线图通过迭代转换，将长为 $k$ 的边序列压缩为单个节点，从而能精确结构化表达具有有限阶局部马尔可夫性质（即 $k$-加性）的复杂路径代价.

\section{CUDA并行框架概述}

为了求解由线图经过规模扩张后的非加性最短路径问题，利用图形处理器（GPU）强大的高并发特性成为提升算法效率的关键.CUDA（Compute Unified Device Architecture）是由 NVIDIA 推出的通用并行计算架构，它通过提供一套统一的软硬件计算框架，使得开发者能够使用 C/C++ 等高级语言直接对 GPU 进行编程，从而将传统的串行复杂图计算任务转化为高度并行的软硬件协同计算.
本节将简要介绍与后续算法设计密切相关的 CUDA 软件架构及其硬件编程模型.

\subsection{CUDA 软件架构}

CUDA 的软件架构在逻辑上基于统一和分层的设计理念，向开发者屏蔽了底层 GPU 硬件的复杂性.其核心思想是“异构计算（Heterogeneous Computing）”，即 CPU 与 GPU 协同工作，各司其职：
\begin{itemize}
    \item \textbf{主机（Host）}：通常指代 CPU 及其内存系统.主机端主要负责处理那些复杂的逻辑控制流、串行计算任务以及任务的调度与资源分配.在最短路径求解任务中，主机负责数据预处理（如线图转换）、内存分配及输出结果的回收.
    \item \textbf{设备（Device）}：通常指代 GPU 及其内存（如全局显存等）.设备端被视作一个高度并行的协处理器，接收来自主机端的数据和计算指令.它主要负责执行拥有极高计算密度和高度并行性的任务，例如大规模图结构中数以万计的节点松弛操作及边权重更新.
\end{itemize}
主机程序通过调用 CUDA 提供的运行时 API（Runtime API）或驱动 API（Driver API），向设备端分发名为“核函数（Kernel）”的并行计算代码，进而触发 GPU 的算力释放.

\subsection{GPU 编程模型与内存结构}

在实际的编程模型中，CUDA 将 GPU 计算抽象为多层级的线程组织架构以及层次化存储结构，以适配网络图结构的不同层次访问需求.

\begin{enumerate}
    \item \textbf{线程组织架构}：CUDA 中的最小执行单元是线程（Thread）.在执行一个 Kernel 时，这些线程被组织起来.若干个线程组成一个线程块（Block），线程块内可以通过共享内存和同步屏障机制实现高效协作；多个并行运算的线程块共同集成一个网络（Grid）.对于最短路径图算法，可以把源点集合或者边集合细粒度地分配给每一个线程或线程块并发处理.
    \item \textbf{层次化内存空间}：为了提升数据吞吐量，CUDA 提供多种存储方式：
    \begin{itemize}
        \item \textbf{全局内存（Global Memory）}：容量最大，主从设备皆可访问，常用于存放整个底层的线图拓扑（如邻接表和权重集）和节点的距离数组.
        \item \textbf{共享内存（Shared Memory）}：位于线程块内部，是一种在Block内的低延迟片上存储介质.在执行松弛迭代时，可将局部更新密集的队列或临时变量放入其中缓冲，大幅降低拉取全局内存的访存延迟.
        \item \textbf{寄存器（Registers）}：每个线程所独占的最快缓存，存储线程执行时的私有上下文变量.
    \end{itemize}
\end{enumerate}

通过上述多线程的阵列调度与多层级的访存优化，基于 CUDA 的并行求解能够实现极高的图计算吞吐率.这不仅能有效抵消“线图拓扑扩张”带来的额外性能开销，更从根本上为现代非加性最短路径问题的高效求解提供了强大的物理算力支撑.

\section{本章小结}

本章先回顾了图的基本概念与加性最短路径的经典定义和算法，再形式化给出非加性最短路径问题及其失效机理，归纳现实应用中的普适性与问题类型；随后总结了改写松弛、线图/状态扩展、动态规划及并行近似等求解思路，并介绍了线图转换要点与CUDA并行框架.通过上述讨论不难发现，现实网络中大量存在的路径依赖型惩罚（如交规转弯限制等），在本质上符合局部马尔可夫性质与有限阶加性（$k$-additive）特征.这也意味着，利用线图及其高阶扩展不仅能有效地捕获此类非加性约束并将其转化为标准的静态模型结构，且高度匹配现代GPU并行的计算范式.正是基于这种理论上的天然契合，下一章将以此为切入点，针对具体的 $k$-加性非加性场景进行更深入的框架设计与相应的并行优化.
